% GENERATED FILE. Edit canonical lessons, then run npm run generate:books.
% canonical-source-hash: fnv1a64:8b7cdfbfda7314af
% canonical-lessons: ES-C580-motivo, ES-C580-a-raiz-de, ES-C580-ya-que, ES-C580-puesto-que, ES-C580-dado-que

\chapter{El porqué}
\label{ch:es-el-porque}
\hlchaptermodality{580}

\begin{chapteropening}
\textbf{By the end of this chapter:} \emph{I can give the reason for something with since and given that.}

Complete the last lesson of chapter 580: I can give the reason for something with since and given that.
\end{chapteropening}

\section[el motivo]{el motivo — the reason, the motive}

\label{lesson:ES-C580-motivo}



\begin{quote}
Before the new one: say the Spanish for to admit, then the Spanish for absurd.
\end{quote}

\subsection*{You'll want to know: el motivo}
\textbf{el motivo} — \textquotedblleft{}the reason, the motive\textquotedblright{}.

\begin{quote}
¿Cuál es el motivo del viaje?
\end{quote}

\begin{quote}
\emph{What is the reason for the trip?}
\end{quote}

\emph{Por este motivo} — for this reason. Close to \emph{la razón}.

\begin{cousinweb}
\emph{motivo} is a learned borrowing from Late Latin \emph{motivus} 'moving', built on \emph{movere} 'to move': a motive is what moves you (EtymDB).
\end{cousinweb}

\begin{grammarlens}[title={What you've built}]
One more word for telling, explaining and describing at B1.
\end{grammarlens}

\subsection*{Your turn}
\begin{itemize}
\raggedright
  \item \emph{Say it:} \emph{el motivo}
  \item \emph{Say it:} \emph{el motivo}, once more
  \item \emph{Say it:} a sentence of your own with \emph{el motivo}
  \item \emph{Recall:} say the Spanish for to admit, then the Spanish for absurd, then say \emph{el motivo} again
\end{itemize}

\begin{tcolorbox}[breakable,colback=teal!4,colframe=teal!35!black,title={Before you move on}]
What does el motivo mean? (\textquotedblleft{}The reason, the motive\textquotedblright{}.) Say it once more.
\end{tcolorbox}
\section[a raíz de]{a raíz de — as a result of, following}

\label{lesson:ES-C580-a-raiz-de}



\begin{quote}
Before the new one: say the Spanish for absurd, then the Spanish for the reason.
\end{quote}

\subsection*{You'll want to know: a raíz de}
\textbf{a raíz de} — \textquotedblleft{}as a result of, following\textquotedblright{}.

\begin{quote}
A raíz del accidente, cambiaron las normas.
\end{quote}

\begin{quote}
\emph{Following the accident, they changed the rules.}
\end{quote}

Followed by a noun; it names the event that started something.

\begin{cousinweb}
\emph{a raíz de} is built on \emph{raíz} 'root', inherited from Latin \emph{radicem}, a form of \emph{radix} 'root' (EtymDB): something grows from the event.
\end{cousinweb}

\begin{grammarlens}[title={What you've built}]
One more word for telling, explaining and describing at B1.
\end{grammarlens}

\subsection*{Your turn}
\begin{itemize}
\raggedright
  \item \emph{Say it:} \emph{a raíz de}
  \item \emph{Say it:} \emph{a raíz de}, once more
  \item \emph{Say it:} a sentence of your own with \emph{a raíz de}
  \item \emph{Recall:} say the Spanish for absurd, then the Spanish for the reason, then say \emph{a raíz de} again
\end{itemize}

\begin{tcolorbox}[breakable,colback=teal!4,colframe=teal!35!black,title={Before you move on}]
What does a raíz de mean? (\textquotedblleft{}As a result of, following\textquotedblright{}.) Say it once more.
\end{tcolorbox}
\section[ya que]{ya que — since, as}

\label{lesson:ES-C580-ya-que}



\begin{quote}
Before the new one: say the Spanish for the reason, then the Spanish for as a result of.
\end{quote}

\subsection*{You'll want to know: ya que}
\textbf{ya que} — \textquotedblleft{}since, as\textquotedblright{}.

\begin{quote}
Ya que estás aquí, ayúdame.
\end{quote}

\begin{quote}
\emph{Since you are here, help me.}
\end{quote}

Gives a reason the listener already knows; it often goes first in the sentence.

\begin{cousinweb}
\emph{ya que} is built on \emph{ya}, inherited from Latin \emph{iam} 'already', and \emph{que}, from Latin \emph{quid} (EtymDB).
\end{cousinweb}

\begin{grammarlens}[title={What you've built}]
One more word for telling, explaining and describing at B1.
\end{grammarlens}

\subsection*{Your turn}
\begin{itemize}
\raggedright
  \item \emph{Say it:} \emph{ya que}
  \item \emph{Say it:} \emph{ya que}, once more
  \item \emph{Say it:} a sentence of your own with \emph{ya que}
  \item \emph{Recall:} say the Spanish for the reason, then the Spanish for as a result of, then say \emph{ya que} again
\end{itemize}

\begin{tcolorbox}[breakable,colback=teal!4,colframe=teal!35!black,title={Before you move on}]
What does ya que mean? (\textquotedblleft{}Since, as\textquotedblright{}.) Say it once more.
\end{tcolorbox}
\section[puesto que]{puesto que — seeing as; since}

\label{lesson:ES-C580-puesto-que}



\begin{quote}
Before the new one: say the Spanish for as a result of, then the Spanish for since.
\end{quote}

\subsection*{You'll want to know: puesto que}
\textbf{puesto que} — \textquotedblleft{}seeing as; since\textquotedblright{}.

\begin{quote}
Puesto que llueve, no salimos.
\end{quote}

\begin{quote}
\emph{Seeing as it is raining, we are not going out.}
\end{quote}

More formal than \emph{porque}; common in writing.

\begin{cousinweb}
\emph{puesto que} is built on \emph{puesto}, the participle of \emph{poner}, which is inherited from Latin \emph{ponere} 'to put' (EtymDB): literally 'it being put that'.
\end{cousinweb}

\begin{grammarlens}[title={What you've built}]
One more word for telling, explaining and describing at B1.
\end{grammarlens}

\subsection*{Your turn}
\begin{itemize}
\raggedright
  \item \emph{Say it:} \emph{puesto que}
  \item \emph{Say it:} \emph{puesto que}, once more
  \item \emph{Say it:} a sentence of your own with \emph{puesto que}
  \item \emph{Recall:} say the Spanish for as a result of, then the Spanish for since, then say \emph{puesto que} again
\end{itemize}

\begin{tcolorbox}[breakable,colback=teal!4,colframe=teal!35!black,title={Before you move on}]
What does puesto que mean? (\textquotedblleft{}Seeing as; since\textquotedblright{}.) Say it once more.
\end{tcolorbox}
\section[dado que]{dado que — given that}

\label{lesson:ES-C580-dado-que}



\begin{quote}
Before the new one: say the Spanish for since, then the Spanish for seeing as.
\end{quote}

\subsection*{You'll want to know: dado que}
\textbf{dado que} — \textquotedblleft{}given that\textquotedblright{}.

\begin{quote}
Dado que no hay dinero, esperamos.
\end{quote}

\begin{quote}
\emph{Given that there is no money, we wait.}
\end{quote}

Formal, like \emph{puesto que}; it is followed by the indicative.

\begin{cousinweb}
\emph{dado que} is built on \emph{dado}, inherited from Latin \emph{datus} 'given', the participle of \emph{dare} 'to give' (EtymDB).
\end{cousinweb}

\begin{grammarlens}[title={What you've built}]
One more word for telling, explaining and describing at B1.
\end{grammarlens}

\subsection*{Your turn}
\begin{itemize}
\raggedright
  \item \emph{Say it:} \emph{dado que}
  \item \emph{Say it:} \emph{dado que}, once more
  \item \emph{Say it:} a sentence of your own with \emph{dado que}
  \item \emph{Recall:} say the Spanish for since, then the Spanish for seeing as, then say \emph{dado que} again
\end{itemize}

\begin{tcolorbox}[breakable,colback=teal!4,colframe=teal!35!black,title={Before you move on}]
What does dado que mean? (\textquotedblleft{}Given that\textquotedblright{}.) Say it once more.
\end{tcolorbox}
